\section{The conductor-\texorpdfstring{$32$}{32} obstruction}
\label{sec:r857-conductor}

\begin{theorem}\label{thm:r857-general}
The ring $R8\_2\_857$ has no categorification over $\C$.
\end{theorem}
\begin{proof}
Reorder the source basis as $(0,1,7,2,5,4,3,6)$ and denote it by
$X_0=\mathbf1,X_1,\ldots,X_7$.  All objects are self-dual.
Multiplication by $X_1$ is the symmetric tridiagonal matrix $J_2$
whose off-diagonal entries are one and whose diagonal is
$(0,1,1,1,1,1,1,2)$.  If $M_j$ denotes multiplication by $X_j$, then
\[
 M_0=I,\qquad M_1=J_2,\qquad
 M_{j+1}=(J_2-I)M_j-M_{j-1}\quad(1\leq j\leq6).
\]
These identities are direct consequences of the source tensor.

Put $\zeta=e^{2\pi i/32}$ and $L=\Q(\zeta+\zeta^{-1})$.
The positive dimension character is
$d_j=\sum_{r=-j}^{j}\zeta^r
=\sin((2j+1)\pi/32)/\sin(\pi/32)$.
Its eight Galois conjugates are distinct and give all ring characters.
Their formal codegrees are the conjugates of
\[
 D=\sum_{j=0}^7d_j^2=4\csc^2(\pi/32),
 \qquad \frac1D=\frac{2-\zeta-\zeta^{-1}}{16}.
\]
For any primitive even Dirichlet character $\chi$ of conductor $32$,
the reciprocal Fourier sum over $\operatorname{Gal}(L/\Q)$ is
\[
 \widehat g(\chi)=-\frac{\tau(\chi)}{16},
 \qquad |\widehat g(\chi)|^2=\frac{32}{16^2}=\frac18.
\]
Here $\tau(\chi)$ is the primitive Gauss sum: summing
$\zeta^a+\zeta^{-a}$ over representatives modulo sign gives precisely
its full sum.  The eight constituents of $I(\mathbf1)$ thus violate
Proposition~\ref{prop:full-I1-bound}: if the spherical dimension field
has index $e$ over $L$, their squared Fourier norm is $e^2$, whereas
the bound is $e/4$.  Hence the ring has no spherical categorification.

Now suppose an arbitrary categorification exists, and take its
canonical sphericalization with sign object $\varepsilon$ of dimension
$-1$.  In the $\varepsilon=1$ sector the eight formal codegrees are
$2\sigma(D)$.  They belong to the spherical dimension field $K$, so
$L\subseteq K$.  Put $e=[K:L]$.  The same eight coordinates in the
full induced-unit compression have squared norm $e^2/4$, with upper
bound $e/4$.  Thus $e=1$ and $[K:\Q]=8$.

Both lifts of every simple are self-dual, since duality preserves its
nonzero spherical dimension while exchanging the lifts negates it.
Choose the positive-second-indicator lifts and pass to the signed
quotient $\varepsilon=-1$.  Its constants satisfy
\[
 |T_{ij}^{k}|\leq N_{ij}^{k},\qquad
 T_{ij}^{k}\equiv N_{ij}^{k}\pmod2,\qquad
 T_{ii}^{j}\equiv N_{ii}^{j}\pmod4.
\]
For the last congruence, a negative-indicator lift can occur in
$X_i^*\otimes X_i$ only with even multiplicity by
\cite[Theorem~3.9]{TPEParity}.  The signed difference is therefore
the original multiplicity minus a multiple of four.
Self-duality gives full symmetry of $T$.  Consequently
$T_{1j}^{j}=T_{jj}^{1}=1$ for $1\leq j\leq6$.

The off-diagonal entries of the signed multiplication matrix of
$X_1$ are $\pm1$.  Successively switching the lifts of
$X_2,\ldots,X_7$ makes them all positive, while leaving the unit,
$X_1$, and all diagonal entries unchanged.  Its matrix is now $J_t$,
with diagonal $(0,1,1,1,1,1,1,t)$ and $t\in\{-2,0,2\}$.
The first seven rows give the same polynomial recurrence for all
signed multiplication matrices.  Thus $\dim X_1$ generates $K$.

If $t=-2$, the bottom-right entry of
$M_2=J_t^2-J_t-I$ is $6$, exceeding the original multiplicity
$N_{27}^{7}=2$.  If $t=0$, then
\[
 \det(xI-J_0)=(x-1)(x+1)(x^2-2x-1)
                    (x^4-4x^3+2x^2+4x-1).
\]
The degree of $\dim X_1$ is at most four, contrary to $[K:\Q]=8$.
Finally, if $t=2$, the recurrence makes every signed multiplication
matrix equal to the original nonnegative one.  The negative lift
multiplicities all vanish, so the eight chosen lifts form a spherical
fusion subcategory equivalent to the original categorification.
This contradicts the spherical exclusion already proved.
\end{proof}

The source-bound certificate verifies all $512$ dimension-character
products in $\Q(\zeta_{32})$, all eight exact Fourier norms, and
the three displayed matrix branches.  The argument is independent
of any construction of the corresponding folded quantum-group datum.
